\chapter{Profile\\นิยามจุดยืนในระบบนิเวศ}

Profile ที่ดีมิใช่ตำแหน่งวิชาการฉบับย่อ แต่เป็นถ้อยแถลงที่ทำให้ผู้อ่านเข้าใจอย่างรวดเร็วว่า Mentee ทำงานกับปัญหาใด ใช้ความเชี่ยวชาญแบบใด และต้องการสร้างคุณค่าให้ใคร การนิยาม Profile จึงเป็นการเลือกจุดยืน มิใช่การรวบรวมทุกสิ่งที่เคยทำ

\learningobjectives{
  \item เขียน Profile ที่เชื่อมความเชี่ยวชาญกับปัญหาและผู้รับคุณค่า
  \item แยกข้อความที่กว้างเกินไปออกจากจุดยืนที่มีความหมาย
  \item ใช้ Profile เป็นเกณฑ์เลือกโอกาสในอนาคต
}

\section{โครงสร้างสามส่วน}
Profile ที่ใช้งานได้ประกอบด้วย \textbf{ความเชี่ยวชาญหลัก} \textbf{ปัญหาที่มุ่งแก้} และ \textbf{คุณค่าหรือกลุ่มผู้ใช้} เช่น \enquote{นักวิจัยด้านเทคโนโลยีหลังการเก็บเกี่ยว ผู้พัฒนาระบบตรวจคุณภาพราคาจับต้องได้ เพื่อช่วยสหกรณ์ลดการสูญเสียและตัดสินใจขายได้แม่นยำขึ้น}

\section{ความลึกและความกว้าง}
การมีความสนใจหลากหลายไม่ใช่ปัญหา แต่ Profile ต้องช่วยให้ผู้อื่นรู้ว่าจะนึกถึง Mentee เมื่อใด โมเดลตัว T มีความลึกหนึ่งด้านและเชื่อมข้ามศาสตร์ได้ ส่วนตัว X สร้างความลึกในหลายด้านที่บรรจบกัน Mentor ควรช่วย Mentee เลือกรูปทรงที่สอดคล้องกับช่วงอาชีพและเป้าหมาย แทนการทำให้ทุกคนเหมือนกัน

\begin{examplebox}
\textbf{ก่อนปรับ:} \enquote{อาจารย์และนักวิจัยด้านเกษตร เทคโนโลยี และนวัตกรรม}\\
\textbf{หลังปรับ:} \enquote{นักวิจัยเกษตรดิจิทัลที่แปลงข้อมูลฟาร์มขนาดเล็กเป็นเครื่องมือวางแผนน้ำและต้นทุน เพื่อเพิ่มความยืดหยุ่นต่อสภาพภูมิอากาศ}\\
ฉบับหลังเปิดพื้นที่ให้ถามต่อเรื่อง Skills, Evidence และ Impact ได้ทันที
\end{examplebox}

\section{ทดสอบ Profile กับผู้ฟังสามกลุ่ม}
ให้ Mentee ทดลองอธิบาย Profile แก่เพื่อนต่างสาขา ผู้ใช้ปลายทาง และผู้สนับสนุนทุน แล้วสังเกตว่าคนแต่ละกลุ่มเข้าใจตรงกันหรือไม่ คำถามที่เกิดซ้ำคือข้อมูลสำคัญสำหรับปรับจุดยืน ความชัดเจนมิได้แปลว่าต้องลดความซับซ้อนของงานทั้งหมด แต่ต้องเลือกประตูเข้าสู่ความซับซ้อนที่เหมาะสม

\diagnostic{เมื่อคนที่ไม่รู้จักท่านอ่าน Profile นี้ เขาควรนึกถึงท่านเมื่อเผชิญปัญหาแบบใด}

\begin{mentornote}
หาก Profile เต็มไปด้วยคำนามนามธรรม เช่น นวัตกรรม ความยั่งยืน หรือความเป็นเลิศ ให้ขอคำกริยาและผู้ได้รับประโยชน์เพิ่ม ความชัดมักเกิดเมื่อเห็นว่า Mentee \enquote{ทำอะไรให้ใคร}
\end{mentornote}

\begin{keytakeaway}
Profile คือสมมติฐานเชิงกลยุทธ์เกี่ยวกับจุดยืนของ Mentee ต้องเฉพาะพอให้คนจดจำ แต่เปิดพอให้เติบโตและสร้างความร่วมมือ
\end{keytakeaway}

\begin{reflection}
เขียน Profile สามฉบับ ความยาว 25 คำ สำหรับผู้ใช้ แหล่งทุน และเพื่อนนักวิจัย เปรียบเทียบว่าคำใดคงที่และคำใดเปลี่ยน จากนั้นเขียนฉบับกลางที่ซื่อตรงและเข้าใจได้โดยทุกกลุ่ม
\blankline\blankline
\end{reflection}

\chaptersummary{Profile แปลประวัติที่ซับซ้อนเป็นจุดยืนที่ผู้อื่นเข้าใจและนำไปเชื่อมโยงได้ จุดยืนที่ชัดช่วยทั้งการสื่อสาร การเลือกโอกาส และการกำหนดหลักฐานที่ควรสร้างต่อไป}
